\setlist[itemize]{noitemsep,topsep=1pt,leftmargin=1.1em}
\setlist[enumerate]{noitemsep,topsep=1pt,leftmargin=1.4em}
\newcommand{\sq}{$\square$\ }

\fbox{\parbox{\dimexpr\linewidth-2\fboxsep-2\fboxrule}{\textbf{КУДА СМОТРЕТЬ (светофор).}\\
\textbf{Нет идеи $>10$ мин} $\to$ раздел «ЗАТУП».\ \ \textbf{WA / RE / TLE / MLE} $\to$ «БАГ» и таблица вердиктов.\\
\textbf{Хочется бросить} $\to$ «Вопросы себе».\ \ \textbf{Не понимаешь, что делать дальше} $\to$ «Ориентиры по времени».}}
\banner{STAY FOCUSED. STAY HARD.}
\subsection{ЗАТУП: нет идеи}
\par \textbf{1. Брут на маленьких.} Посмотри на ответы глазами: есть ли паттерн? бывает ли вообще \texttt{-1}? сколько решений? Гипотезу проверяй брутом, а не верой.

\par \fbox{\parbox{\dimexpr\linewidth-2\fboxsep-2\fboxrule}{\textbf{4 СТОП-ВОПРОСА} -- в эти моменты не иди дальше, а спроси себя.\\
\textbf{A. Нашёл правило, сработало на примерах:} «ПОЧЕМУ в конце больше ничего нельзя сделать?» Не можешь ответить -- правило неполное. Ищи контрпример руками на 3--4 элементах (лесенка: 1-2, 2-3, 1-2, 2-3).\\
\textbf{B. Застрял 10--15 минут:} «Какие ЕЩЁ объекты есть в задаче, кроме тех, что в условии?» Выпиши ВСЕ: элементы, точки, позиции, пары, границы, вклады, разности. Потом: «СКОЛЬКО их участвует в одном шаге и что ровно меняется?» (лесенка: не перекладины, а точки; на перекладине ровно две).\\
\textbf{C. Хочешь ускорить:} «Откуда квадрат -- произведение ЧЕГО на ЧТО?» Выпиши все множители (число стартов $\times$ стоимость старта и т.д.); для каждого: можно ли уменьшить? Два пути: делать меньше (стартов/состояний) или переиспользовать работу.\\
\textbf{D. Доказал лемму:} «Что она ЗАПРЕЩАЕТ? Какие слова теперь нельзя использовать?» И после каждого хода: «Что изменилось? Появилось ли новое?»}}
\par \textbf{2. Что НЕ зависит} от выбора жюри / разбиения / порядка? Что стоит на фиксированных местах? (например, границы цифр на фиксированных позициях.)

\par \textbf{3. Смени взгляд:} не числа, а столбцы цифр; не разности, а абсолютные позиции; не элемент, а его вклад; локальные случаи $\to$ вся строка.
\par \textbf{4. Сколько РЕАЛЬНО состояний?} (делителей $\le10^5$, мемо по остатку). Раздели на маленькие и большие (большое бывает максимум один раз).
\par \textbf{5. Переформулируй сумму:} по перестановкам $\to$ по подмножествам с весом $k!(m-k)!$; среднее/сумма $\to$ по позициям отдельно (линейность); «ровно $d$» $=S(\le d)-S(\le d-1)$, где $S$ считается просто; ответ $=$ всё $-$ дополнение.
\par \textbf{6. Кусочно-многочленная функция} $\to$ сумму на куске интерполируй (Лагранж, значения в $0..\deg$).
\par \textbf{7. Разбираешь кучу случаев} $\to$ скорее всего, не тот взгляд. Вернись к B и к пункту 2.


\banner{NO ONE IS COMING TO SAVE YOU. SOLVE IT.}
\subsection{Хочется сдаться? Вопросы себе}
\par\begin{enumerate}
\item Что я чувствую: \emph{усталость} или \emph{реальный предел}? Мозг сдаётся, когда ты на 40\%. (Правило 40\%.)
\item Какой самый маленький шаг я могу сделать \emph{прямо сейчас}, за 5 минут? Сделай его.
\item Что я могу контролировать? Не результат, а: перечитать условие, написать брут, пройти чек-лист, сказать идею вслух.
\item Какую задачу мы ещё \emph{не пробовали}? Сдача не только на одной.
\item Что я посоветовал бы напарнику, если бы он сидел вот так?
\item Через час я буду жалеть, что остановился, или что продолжил?
\item Если бы это был последний час турнира -- что бы я делал? Делай это уже сейчас.
\end{enumerate}
\par \textbf{Мантры.} Не жди мотивации -- мотивация приходит и уходит, а дисциплина остаётся. Один раз сдашься -- мозг запомнит. Не сдашься -- запомнит это. Сделай ещё один шаг, ещё один тест, ещё одну попытку.
\banner{WHO'S GONNA CARRY THE BOATS AND THE LOGS?}
\banner{WE ARE, MOTHERFU***RS. WE ARE.}
